\section{Problemet med en byte-ström}
\label{sec:stream}

Två system kopplas ihop med en seriell länk --- UART, RS-485, SPI, vad som helst. Det ena skickar
bytes, det andra tar emot dem. Det är allt länken gör.

Anta att den sändande noden vill skicka temperaturen 23 grader och luftfuktigheten 45 procent.
Den skickar två bytes:

\begin{console}
17 2D
\end{console}

Mottagaren får två bytes. Och nu kommer frågorna, i tur och ordning:

\begin{itemize}
  \item \textbf{Var börjar meddelandet?} Om mottagaren kopplades in mitt i en överföring är
    \code{17} kanske den andra byten i något annat.
  \item \textbf{Hur långt är det?} Kommer det fler bytes, eller var det allt?
  \item \textbf{Vad betyder de?} Att \code{17} är temperaturen och \code{2D} fukten är en
    överenskommelse som inte finns någonstans i strömmen.
  \item \textbf{Kom de fram hela?} En störning som ändrade \code{17} till \code{16} ser ut precis
    som en temperatur på 22 grader.
  \item \textbf{Vem skickade dem, och till vem?} På en buss med fler än två noder är det inte
    självklart.
\end{itemize}

Ingen av frågorna går att besvara ur bytesen själva. Svaren måste läggas till, och det som lägger
till dem är ett \textbf{protokoll}.

\subsection{Frame kontra ström}

En \textbf{frame} är en avgränsad, självbeskrivande enhet i strömmen. Den säger var den börjar,
hur lång den är, vad den innehåller, vem den kommer från, vem den är till, och om den kom fram
hel.

Strömmen är fortfarande bara bytes --- det är den alltid. Skillnaden är att bytesen nu är ordnade
efter en överenskommelse som båda sidor känner till, och som går att kontrollera.

\begin{aside}
\note Det här är den enda idén i kapitlet, och resten av boken är konsekvenser av den. Både det
egna protokollet i \kapref{ch:parsning} och CAN i \kapref{ch:can} löser samma fem frågor. De
skiljer sig i \emph{var} lösningen ligger: i din kod, eller i kisel.
\end{aside}
